\chapter{System Design}

\section{System Architecture Overview}

The forecasting system follows a modular, pipeline-based architecture organized into five distinct stages: data ingestion, preprocessing, model training and optimization, evaluation, and deployment. Figure~\ref{fig:arch_overview} illustrates the overall system architecture.

\begin{figure}[H]
\centering
\fbox{\begin{minipage}{0.9\textwidth}
\centering
\vspace{0.2in}
\textbf{Data Ingestion} $\rightarrow$ \textbf{Preprocessing} $\rightarrow$ \textbf{Model Training} $\rightarrow$ \textbf{Evaluation} $\rightarrow$ \textbf{Deployment}
\vspace{0.1in}

Open-Meteo API $\rightarrow$ Raw CSV $\rightarrow$ Normalized Tensors $\rightarrow$ 14 Trained Models $\rightarrow$ ONNX + Pi
\vspace{0.2in}
\end{minipage}}
\caption{High-Level System Architecture}
\label{fig:arch_overview}
\end{figure}

\section{Data Pipeline Design}

\subsection{Data Ingestion}

The data ingestion module interfaces with the Open-Meteo Air Quality API (\texttt{https://air-quality-api.open-meteo.com/v1/air-quality}) to fetch hourly observations. The API returns JSON responses containing the specified pollutant concentrations and meteorological parameters for the requested geographic coordinates (Hyderabad: 17.3850\textdegree N, 78.4867\textdegree E).

Key design decisions include:
\begin{itemize}
    \item \textbf{Minimal API calls}: One API call retrieves the full year of historical data, minimizing network overhead.
    \item \textbf{Structured storage}: Data is cached locally as CSV files for reproducibility and offline access.
    \item \textbf{Timestamp handling}: All timestamps are converted to UTC and indexed for consistent temporal operations.
\end{itemize}

\subsection{Data Preprocessing}

The preprocessing pipeline performs the following operations in sequence:

\begin{enumerate}
    \item \textbf{Missing Value Imputation}: Gaps in the time series are identified and filled using linear interpolation. This preserves temporal continuity without introducing artifacts.
    \item \textbf{Feature Selection}: Seven pollutant features (PM$_{2.5}$, PM$_{10}$, CO, NO$_2$, SO$_2$, O$_3$, European AQI) are selected as input features. PM$_{2.5}$ is designated as the target variable.
    \item \textbf{Normalization}: Min-Max scaling is applied to all features, transforming them to the [0, 1] range:
    \begin{equation}
        x_{\text{scaled}} = \frac{x - x_{\min}}{x_{\max} - x_{\min}}
    \end{equation}
    \item \textbf{Sliding Window Construction}: Input-output pairs are created using a 168-hour sliding window (7 days of history) to predict the next 24 hours. For a time series of length $T$, this yields $T - 168 - 24$ training samples.
\end{enumerate}

\subsection{Data Splitting Strategy}

The dataset is split chronologically rather than randomly to respect temporal ordering and prevent look-ahead bias:
\begin{itemize}
    \item Training: 70\% (samples 0--6,149)
    \item Validation: 15\% (samples 6,150--7,466)
    \item Test: 15\% (samples 7,467--8,784)
\end{itemize}

\section{Model Architectures}

This section presents the detailed architecture of each of the 14 forecasting models evaluated in this study. All deep learning models share a common input-output interface: input tensors of shape $(B, 168, 7)$ where $B$ is batch size, 168 is the lookback window, and 7 is the number of features; output tensors of shape $(B, 24)$ representing the 24-hour PM$_{2.5}$ forecast.

\subsection{ARIMA (AutoRegressive Integrated Moving Average)}

The ARIMA model, denoted as ARIMA$(p, d, q)$, is defined by three parameters: $p$ (autoregressive order), $d$ (degree of differencing), and $q$ (moving average order). The model can be expressed as:

\begin{equation}
    \phi(B)(1 - B)^d y_t = \theta(B)\varepsilon_t
\end{equation}

where $B$ is the backshift operator, $\phi(B) = 1 - \phi_1 B - \cdots - \phi_p B^p$ is the AR polynomial, $\theta(B) = 1 + \theta_1 B + \cdots + \theta_q B^q$ is the MA polynomial, and $\varepsilon_t \sim \mathcal{N}(0, \sigma^2)$ is white noise.

For our implementation, we selected ARIMA(5, 1, 2) after grid search optimization. ARIMA is a univariate model, using only the historical PM$_{2.5}$ values for prediction, ignoring other pollutant features.

\subsection{SARIMA (Seasonal ARIMA)}

SARIMA extends ARIMA by incorporating seasonal components, denoted as SARIMA$(p, d, q)(P, D, Q)_s$:

\begin{equation}
    \phi(B)\Phi(B^s)(1 - B)^d(1 - B^s)^D y_t = \theta(B)\Theta(B^s)\varepsilon_t
\end{equation}

where $\Phi(B^s)$ and $\Theta(B^s)$ are the seasonal AR and MA polynomials, $s$ is the seasonal period, and $P, D, Q$ are the seasonal orders. We used SARIMA(2, 1, 1)(1, 1, 1)$_{24}$, with a 24-hour seasonal period corresponding to the diurnal cycle.

\subsection{VAR (Vector Autoregression)}

VAR models the joint dynamics of multiple time series. A VAR($p$) model with $k$ variables is expressed as:

\begin{equation}
    \mathbf{y}_t = \mathbf{c} + \mathbf{A}_1 \mathbf{y}_{t-1} + \cdots + \mathbf{A}_p \mathbf{y}_{t-p} + \boldsymbol{\varepsilon}_t
\end{equation}

where $\mathbf{y}_t \in \mathbb{R}^k$, $\mathbf{c}$ is a constant vector, each $\mathbf{A}_i \in \mathbb{R}^{k \times k}$ is a coefficient matrix, and $\boldsymbol{\varepsilon}_t \sim \mathcal{N}(\mathbf{0}, \boldsymbol{\Sigma})$. We used VAR(2) with $k = 7$ variables (all pollutant features plus target).

\subsection{Simple RNN}

The vanilla RNN processes sequential data through a recurrent hidden state:

\begin{equation}
    \mathbf{h}_t = \tanh(\mathbf{W}_{ih} \mathbf{x}_t + \mathbf{b}_{ih} + \mathbf{W}_{hh} \mathbf{h}_{t-1} + \mathbf{b}_{hh})
\end{equation}

Our RNN implementation uses 2 recurrent layers with 128 hidden units each and dropout of 0.2 between layers. The final hidden state is passed through a linear layer to produce the 24-hour forecast.

\subsection{LSTM (Long Short-Term Memory)}

LSTM addresses the vanishing gradient problem through a gating mechanism with three gates:

\begin{align}
    \mathbf{f}_t &= \sigma(\mathbf{W}_f \cdot [\mathbf{h}_{t-1}, \mathbf{x}_t] + \mathbf{b}_f) \quad \text{(forget gate)} \\
    \mathbf{i}_t &= \sigma(\mathbf{W}_i \cdot [\mathbf{h}_{t-1}, \mathbf{x}_t] + \mathbf{b}_i) \quad \text{(input gate)} \\
    \mathbf{o}_t &= \sigma(\mathbf{W}_o \cdot [\mathbf{h}_{t-1}, \mathbf{x}_t] + \mathbf{b}_o) \quad \text{(output gate)} \\
    \tilde{\mathbf{c}}_t &= \tanh(\mathbf{W}_c \cdot [\mathbf{h}_{t-1}, \mathbf{x}_t] + \mathbf{b}_c) \quad \text{(candidate cell)} \\
    \mathbf{c}_t &= \mathbf{f}_t \odot \mathbf{c}_{t-1} + \mathbf{i}_t \odot \tilde{\mathbf{c}}_t \quad \text{(cell state)} \\
    \mathbf{h}_t &= \mathbf{o}_t \odot \tanh(\mathbf{c}_t) \quad \text{(hidden state)}
\end{align}

where $\sigma$ is the sigmoid activation and $\odot$ denotes element-wise multiplication. Our LSTM uses 2 layers, 128 hidden units, and dropout of 0.2. The model has 147,073 trainable parameters.

\subsection{GRU (Gated Recurrent Unit)}

GRU simplifies the LSTM by merging the cell state and hidden state, and combining the forget and input gates into a single update gate:

\begin{align}
    \mathbf{z}_t &= \sigma(\mathbf{W}_z \cdot [\mathbf{h}_{t-1}, \mathbf{x}_t] + \mathbf{b}_z) \quad \text{(update gate)} \\
    \mathbf{r}_t &= \sigma(\mathbf{W}_r \cdot [\mathbf{h}_{t-1}, \mathbf{x}_t] + \mathbf{b}_r) \quad \text{(reset gate)} \\
    \tilde{\mathbf{h}}_t &= \tanh(\mathbf{W}_h \cdot [\mathbf{r}_t \odot \mathbf{h}_{t-1}, \mathbf{x}_t] + \mathbf{b}_h) \quad \text{(candidate hidden)} \\
    \mathbf{h}_t &= (1 - \mathbf{z}_t) \odot \mathbf{h}_{t-1} + \mathbf{z}_t \odot \tilde{\mathbf{h}}_t
\end{align}

Our GRU implementation uses 2 layers, 128 hidden units, and dropout of 0.2, resulting in 110,593 parameters --- 25\% fewer than the equivalent LSTM.

\subsection{BiLSTM (Bidirectional LSTM)}

BiLSTM processes the input sequence in both forward and backward directions, concatenating the hidden states:

\begin{equation}
    \mathbf{h}_t^{\text{bi}} = [\overrightarrow{\mathbf{h}}_t ; \overleftarrow{\mathbf{h}}_t]
\end{equation}

This doubles the effective hidden dimension, enabling the model to capture context from both past and future time steps within the lookback window. Our BiLSTM uses 2 bidirectional LSTM layers with 128 hidden units per direction, resulting in 491,521 parameters.

\subsection{CNN-LSTM}

The CNN-LSTM architecture combines 1D convolutional feature extraction with LSTM temporal modeling:

\begin{enumerate}
    \item \textbf{Convolutional Block}: Two 1D convolutional layers with 64 and 128 filters (kernel size 3), each followed by ReLU activation and dropout (0.2). This extracts local temporal patterns from the input window.
    \item \textbf{LSTM Block}: Two LSTM layers with 128 hidden units process the convolved features, capturing long-range dependencies.
    \item \textbf{Output Layer}: A fully connected layer maps the final LSTM hidden state to 24 output predictions.
\end{enumerate}

The convolutional front-end effectively reduces the sequence length while preserving salient features, making the subsequent LSTM processing more efficient.

\subsection{CNN-GRU}

CNN-GRU follows the same architectural pattern as CNN-LSTM but substitutes GRU layers for the recurrent component. This provides approximately 25\% fewer recurrent parameters while maintaining comparable performance.

\subsection{BiLSTM with Attention}

This architecture augments the BiLSTM with an additive attention mechanism that learns to weight the importance of each time step:

\begin{align}
    e_t &= \mathbf{v}^T \tanh(\mathbf{W}_a \mathbf{h}_t^{\text{bi}} + \mathbf{b}_a) \\
    \alpha_t &= \frac{\exp(e_t)}{\sum_{j=1}^{T} \exp(e_j)} \\
    \mathbf{c} &= \sum_{t=1}^{T} \alpha_t \mathbf{h}_t^{\text{bi}}
\end{align}

where $\mathbf{v}$, $\mathbf{W}_a$, and $\mathbf{b}_a$ are learnable attention parameters, $\alpha_t$ are the attention weights, and $\mathbf{c}$ is the context vector fed to the output layer. The attention mechanism enables the model to focus on the most informative time steps for prediction.

\subsection{Transformer}

The Transformer architecture replaces recurrence with multi-head self-attention:

\begin{equation}
    \text{Attention}(Q, K, V) = \text{softmax}\left(\frac{QK^T}{\sqrt{d_k}}\right) V
\end{equation}

\begin{equation}
    \text{MultiHead}(Q, K, V) = \text{Concat}(\text{head}_1, \ldots, \text{head}_h) W^O
\end{equation}

Our Transformer implementation uses 4 encoder layers with 4 attention heads, a model dimension ($d_{\text{model}}$) of 128, feed-forward dimension of 512, and dropout of 0.1. Positional encoding is added to the input embeddings to provide temporal ordering information, as the self-attention mechanism is inherently permutation-invariant.

\subsection{Informer}

The Informer architecture introduces three key innovations for efficient long-sequence forecasting:

\begin{enumerate}
    \item \textbf{ProbSparse Self-Attention}: Instead of computing attention for all query-key pairs, ProbSparse attention selects only the top-$u$ most influential queries based on a sparsity measurement:
    \begin{equation}
        M(q_i, K) = \max_j \left\{\frac{q_i k_j^T}{\sqrt{d}}\right\} - \frac{1}{L_K} \sum_{j=1}^{L_K} \frac{q_i k_j^T}{\sqrt{d}}
    \end{equation}
    
    \item \textbf{Self-Attention Distilling}: Between encoder layers, a distilling operation with 1D max-pooling progressively reduces the sequence length, prioritizing dominant attention features.
    
    \item \textbf{Generative Decoder}: The decoder produces the entire forecast sequence in a single forward pass, rather than auto-regressively.
\end{enumerate}

Our Informer uses 3 encoder layers, 2 decoder layers, $d_{\text{model}} = 128$, and 4 attention heads.

\subsection{Autoformer}

Autoformer introduces an Auto-Correlation mechanism as a drop-in replacement for self-attention:

\begin{equation}
    \mathcal{R}_{XX}(\tau) = \lim_{L \to \infty} \frac{1}{L} \sum_{t=1}^{L} X_t X_{t - \tau}
\end{equation}

The architecture explicitly decomposes the time series into trend and seasonal components through a series of moving average operations within each Auto-Correlation block:

\begin{align}
    X_t^{\text{trend}} &= \text{AvgPool}(\text{Padding}(X)) \\
    X_t^{\text{seasonal}} &= X - X_t^{\text{trend}}
\end{align}

This progressive decomposition provides a strong inductive bias for seasonal time series, which is especially relevant for air quality data with strong diurnal and weekly cycles. Our Autoformer uses 2 encoder layers, 1 decoder layer, and $d_{\text{model}} = 128$.

\subsection{ST-GCN (Spatio-Temporal Graph Convolutional Network)}

ST-GCN models the multivariate time series as a graph where each pollutant is a node. The adjacency matrix is learned from the data using the spatial attention mechanism:

\begin{equation}
    \mathbf{A} = \text{softmax}(\text{ReLU}(\mathbf{E}_1 \mathbf{E}_2^T))
\end{equation}

where $\mathbf{E}_1$ and $\mathbf{E}_2$ are learnable node embeddings. Graph convolutions propagate information along the learned graph edges:

\begin{equation}
    \mathbf{H}^{(l+1)} = \sigma(\tilde{\mathbf{D}}^{-\frac{1}{2}} \tilde{\mathbf{A}} \tilde{\mathbf{D}}^{-\frac{1}{2}} \mathbf{H}^{(l)} \mathbf{W}^{(l)})
\end{equation}

Temporal convolutions with gated activation capture sequential patterns. Our ST-GCN uses 2 spatio-temporal blocks with 64 hidden channels and a Chebyshev polynomial order of 3 for graph convolutions.

\section{Hyperparameter Optimization}

\subsection{Optuna Framework}

Hyperparameter optimization is conducted using Optuna, a Bayesian optimization framework that uses the Tree-structured Parzen Estimator (TPE) algorithm. TPE models $p(x|y)$ and $p(y)$ separately, where $x$ represents hyperparameters and $y$ represents the objective value (validation RMSE):

\begin{equation}
    p(x|y) = \begin{cases}
        l(x) & \text{if } y < y^* \\
        g(x) & \text{if } y \geq y^*
    \end{cases}
\end{equation}

The algorithm selects hyperparameters that maximize the ratio $l(x)/g(x)$, effectively sampling from regions of the hyperparameter space likely to produce good results.

\subsection{Search Space}

Table~\ref{tab:search_space} presents the hyperparameter search space for each model family.

\begin{table}[H]
\centering
\caption{Hyperparameter Search Space by Model Family}
\label{tab:search_space}
\begin{tabular}{p{3cm}p{2.8cm}p{6.5cm}}
\toprule
\textbf{Family} & \textbf{Parameter} & \textbf{Range} \\
\midrule
\multirow{3}{*}{Statistical} & ARIMA $p$ & $\{1, \ldots, 10\}$ \\
 & ARIMA $q$ & $\{0, \ldots, 5\}$ \\
 & VAR order & $\{1, \ldots, 5\}$ \\
\midrule
\multirow{4}{*}{RNN/LSTM/GRU} & Hidden size & $\{64, 128, 256\}$ \\
 & Num layers & $\{1, 2, 3\}$ \\
 & Dropout & $[0.1, 0.5]$ \\
 & Learning rate & Log-uniform $[10^{-4}, 10^{-2}]$ \\
\midrule
\multirow{5}{*}{CNN-RNN Hybrids} & CNN filters & $\{32, 64, 128\}$ \\
 & Kernel size & $\{3, 5, 7\}$ \\
 & RNN hidden & $\{64, 128\}$ \\
 & RNN layers & $\{1, 2\}$ \\
 & Learning rate & Log-uniform $[10^{-4}, 10^{-2}]$ \\
\midrule
\multirow{5}{*}{Transformers} & $d_{\text{model}}$ & $\{64, 128, 256\}$ \\
 & Num heads & $\{4, 8\}$ \\
 & Encoder layers & $\{2, 3, 4\}$ \\
 & FF dimension & $\{256, 512\}$ \\
 & Learning rate & Log-uniform $[10^{-4}, 10^{-2}]$ \\
\bottomrule
\end{tabular}
\end{table}

\subsection{Optimization Budget}

Each model undergoes 20 Optuna trials. The median wall-clock time for 20 trials ranges from 5 minutes (RNN) to 4.2 hours (BiLSTM with Attention). The objective function minimizes validation set RMSE.

\section{Training Configuration}

All deep learning models share a common training configuration:

\begin{itemize}
    \item \textbf{Loss Function}: Mean Squared Error (MSE)
    \item \textbf{Optimizer}: Adam ($\beta_1 = 0.9$, $\beta_2 = 0.999$)
    \item \textbf{Learning Rate Scheduler}: ReduceLROnPlateau with factor 0.5 and patience of 10 epochs
    \item \textbf{Early Stopping}: Patience of 20 epochs on validation loss, restoring best weights
    \item \textbf{Maximum Epochs}: 200
    \item \textbf{Batch Size}: 64
    \item \textbf{Random Seed}: 42 (ensures reproducibility)
\end{itemize}

\section{Deployment Architecture}

The deployment pipeline converts trained PyTorch models to ONNX format for efficient inference on the Raspberry Pi 4:

\begin{enumerate}
    \item \textbf{Model Export}: Traced PyTorch models are exported using \texttt{torch.onnx.export()} with dynamic axes for batch size.
    \item \textbf{ONNX Validation}: Exported models are validated using \texttt{onnx.checker.check\_model()}.
    \item \textbf{Quantization}: Optional INT8 dynamic quantization reduces model size and inference latency.
    \item \textbf{Inference Session}: ONNX Runtime creates an inference session with CPU execution provider optimized for ARM64.
\end{enumerate}

The Raspberry Pi 4 deployment setup includes:
\begin{itemize}
    \item \textbf{OS}: Raspberry Pi OS (64-bit)
    \item \textbf{Runtime}: ONNX Runtime 1.16 with ARM64 optimizations
    \item \textbf{Storage}: All 14 ONNX models stored on microSD
    \item \textbf{Benchmarking}: Per-model latency measured over 100 inference runs with warmup iterations
\end{itemize}
